\section{Prompt and Judge Details}
\subsection{Prompt contract}
Each run uses a two-message contract:
\begin{itemize}
  \item \textbf{System}: ``Solve the given prompt exactly. Return ONLY code (no explanation, no markdown). Follow alias mapping strictly.''
  \item \textbf{User}: versioned prompt text (\texttt{docs/archive/prompt-versions/v*.md}) containing alias mapping and task statement.
\end{itemize}

\subsection{Input prompt structure}
Versioned prompts include:
\begin{enumerate}
  \item alias map block (e.g., token\(\rightarrow\)alias pairs),
  \item task statement (Fibonacci-style function + loop output requirement),
  \item strict-format instruction.
\end{enumerate}

\subsection{Deterministic judge protocol}
The judge in \texttt{scripts/run\_gpt54\_eval.py} applies the following checks:
\begin{enumerate}
  \item Alias compliance: required alias presence and original keyword leakage detection.
  \item Normalization validity: alias\(\rightarrow\)Python normalization then AST parseability.
  \item Skeleton checks:
  \begin{itemize}
    \item function \texttt{fib} existence,
    \item base case \texttt{if n <= 1: return n},
    \item recursive return \texttt{fib(n-1) + fib(n-2)},
    \item loop \texttt{for i in range(7): print(i, fib(i))}.
  \end{itemize}
\end{enumerate}

\subsection{Pass/fail and score}
Pass requires zero violations. Score starts at 100 and decreases by fixed penalties per violation. This provides a strict and reproducible criterion for cross-model comparison under transformed syntax constraints.

\subsection{Failure taxonomy source}
Violation classes are directly emitted into result JSON and aggregated in summary artifacts under \texttt{docs/research/results/}. This enables exact claim-to-artifact traceability.
